\chapter*{กิตติกรรมประกาศ}
\addcontentsline{toc}{chapter}{กิตติกรรมประกาศ}

การจัดทำตำราวิชาการและชุดเอกสารหลักสูตรเล่มนี้ สำเร็จลุล่วงได้อย่างสมบูรณ์ด้วยความร่วมมือ ร่วมแรง และร่วมใจจากคณาจารย์ นักวิจัย และบุคลากรของคณะวิทยาศาสตร์และเทคโนโลยี มหาวิทยาลัยราชภัฏรำไพพรรณี ที่ได้ร่วมกันถ่ายทอดองค์ความรู้และประสบการณ์วิจัยเชิงประยุกต์ตลอดระยะเวลากว่าทศวรรษ

คณะผู้เขียนขอกราบขอบพระคุณคณะผู้บริหารมหาวิทยาลัยราชภัฏรำไพพรรณี และคณบดีคณะวิทยาศาสตร์และเทคโนโลยี ที่ได้ให้การสนับสนุนด้านนโยบาย งบประมาณ และสิ่งอำนวยความสะดวกในการวิจัยพัฒนา รวมถึงห้องปฏิบัติการวิทยาศาสตร์และฟาร์มทดสอบต้นแบบ

ขอขอบพระคุณคณะวิทยากรผู้ทรงคุณวุฒิทุกท่าน ได้แก่ ผู้ช่วยศาสตราจารย์ อรรถกร คำฉัตร, ผู้ช่วยศาสตราจารย์ ดร.ณมนรัก คำฉัตร, รองศาสตราจารย์ ดร.นิภัทร เปี่ยมอรุณ, ผู้ช่วยศาสตราจารย์ ดร.จิรภัทร จันทมาลี, ผู้ช่วยศาสตราจารย์ นิคม รัตนโรจนกุล, รองศาสตราจารย์ ดร.นันทพร มูลรังษี และผู้ช่วยศาสตราจารย์ ดร.สุพัตรา รักษาพรต ที่ได้สละเวลาทุ่มเทจัดเตรียมเนื้อหาวิชาการ ตรวจสอบความถูกต้องทางวิทยาศาสตร์ และพัฒนาระบบซอฟต์แวร์ต้นแบบภาคสนาม

ขอขอบพระคุณเครือข่ายแปลงวิจัยและชาวสวนผลไม้ในจังหวัดจันทบุรี ตราด และระยอง รวมถึงผู้ประกอบการฟาร์มเพาะเลี้ยงกุ้งขาวและสัตว์น้ำชายฝั่ง ที่ได้เอื้อเฟื้อพื้นที่ในการเก็บตัวอย่างดิน สภาพแวดล้อมทางชลศาสตร์ และการทดสอบระบบสมองกลฝังตัวหน้างานจริง ซึ่งทำให้ข้อมูลและผลการทดลองในตำราเล่มนี้มีความน่าเชื่อถือและสอดคล้องกับสภาพปัญหาที่แท้จริง

คุณงามความดีและประโยชน์อันพึงมีจากตำราเล่มนี้ คณะผู้เขียนขอน้อมอุทิศแด่บูรพาจารย์ ผู้มีพระคุณ และสถาบันการศึกษาทุกแห่งที่ได้ประสิทธิ์ประสาทวิชาความรู้ให้แก่คณะผู้เขียนเสมอมา

\vspace{1.5cm}
\begin{flushright}
    \textbf{คณะผู้จัดทำหลักสูตร AIoT 2570}\\
    มหาวิทยาลัยราชภัฏรำไพพรรณี
\end{flushright}
\clearpage
